\documentclass[11pt, a4paper]{article}
\usepackage[a4paper, margin=2.5cm]{geometry}
\usepackage[utf8]{inputenc}
\usepackage[T1]{fontenc}
\usepackage[spanish, es-tabla]{babel}
\usepackage{booktabs}
\usepackage{tabularx}
\usepackage{longtable}
\usepackage{enumitem}
\usepackage{xcolor}
\usepackage{graphicx}
\usepackage{hyperref}

\hypersetup{
    colorlinks=true,
    linkcolor=blue!80!black,
    urlcolor=blue!80!black,
    citecolor=blue!80!black
}

\setlist[itemize]{label=-}

\title{\textbf{Trabajo Individual: ISO/IEC/IEEE 29148:2018 y Tipos de Requisitos de un Sistema Conocido (Mercado Libre)}\\ \large Ingeniería de Requerimientos (ISR-401) -- 4.\textsuperscript{o} Nivel Paralelo A}
\author{\textbf{Estudiante:} Yerai David Chalacan Toala \\ \textbf{Docente:} Gleiston Cicerón Guerrero Ulloa, PhD \\ \texttt{gguerrero@uteq.edu.ec}}
\date{Semana 2 (2 de noviembre - 8 de noviembre de 2026)}

\begin{document}

\maketitle

\section{Introducción y Descripción del Sistema}
En esta actividad académica analizo y aplico los conceptos principales de la norma internacional ISO/IEC/IEEE 29148:2018 \cite{iso29148} junto con el modelo de calidad ISO/IEC 25010:2023 \cite{iso25010}. Para llevar los conceptos teóricos a la práctica, elegí la aplicación móvil de comercio electrónico \textbf{Mercado Libre} (versión Ecuador), una plataforma que utilizo habitualmente para mis compras.

Mercado Libre funciona como un marketplace digital interactivo enfocado en conectar a compradores y vendedores de productos nuevos o usados. Entre los usuarios que participan en la plataforma están los compradores particulares, vendedores individuales, tiendas oficiales y las empresas de logística encargadas de los envíos. En mi experiencia diaria de uso, las funciones principales abarcan desde buscar productos y aplicar filtros, hasta revisar el stock disponible, procesar el pago con tarjetas de débito o crédito y hacer el seguimiento de la entrega del paquete hasta el domicilio.

\section{Resumen Analítico de la Norma ISO/IEC/IEEE 29148:2018}

\subsection{Construcción de un Requisito Bien Formado (Cláusula 5.2.4)}
Al revisar la cláusula 5.2.4 de la norma ISO/IEC/IEEE 29148:2018 \cite{iso29148}, observo que la redacción de un requisito no debe hacerse de forma libre o informal, sino siguiendo una sintaxis bien definida para evitar malentendidos entre desarrolladores y evaluadores. La norma propone que todo requisito bien estructurado cuente con tres partes básicas: un sujeto claramente identificado, una condición o contexto en el que opera, y la acción específica expresada como un resultado medible.

En cuanto a las reglas gramaticales, la norma enfatiza el uso del verbo obligatorio \textit{«shall»} (traducido normalmente como «deberá» o «el sistema deberá») para marcar funciones indispensables. Cuando se trata de eventos desencadenados por el usuario o por el sistema, se utiliza la fórmula: \textit{«Cuando [ocurra cierta condición], el [sistema/sujeto] deberá [realizar la acción]»}. Es fundamental diferenciar entre \textit{«shall»} (obligación normativa), \textit{«should»} (sugerencia no obligatoria) y \textit{«may»} (opción permitida), descartando palabras vagas que impidan probar el requisito.

\subsection{Características de un Requisito Individual (Cláusula 5.2.5)}
La cláusula 5.2.5 \cite{iso29148} establece nueve propiedades indispensables que debe tener cada requerimiento para garantizar que esté bien redactado antes de pasar a la fase de diseño:
\begin{enumerate}
    \item \textbf{Necesario (Necessary):} Refleja una necesidad real y valiosa para los usuarios o el negocio, sin la cual el sistema estaría incompleto.
    \item \textbf{Adecuado al diseño (Implementation free):} Explica claramente qué debe hacer el software sin imponer una solución de código o tecnología específica.
    \item \textbf{Inambiguo (Unambiguous):} Su redacción es tan clara que solo admite una única interpretación posible para cualquier lector.
    \item \textbf{Completo (Complete):} Contiene toda la información y condiciones requeridas en una sola declaración, sin dejar aspectos a la imaginación.
    \item \textbf{Singular (Singular):} Describe una sola función o restricción por enunciado, evitando mezclar varias ideas en una misma frase.
    \item \textbf{Factible (Feasible):} Es posible de construir dentro de los plazos, presupuesto y limitaciones técnicas reales del proyecto.
    \item \textbf{Verificable (Verifiable):} Permite comprobar objetivamente si se cumplió o no mediante pruebas, inspecciones o demostraciones.
    \item \textbf{Correcto (Correct):} Coincide exactamente con la necesidad manifestada por el usuario o interesado.
    \item \textbf{Conforme (Conforming):} Respeta la estructura, el estilo y la sintaxis definidos por la norma del proyecto.
\end{enumerate}

\subsection{Características de un Conjunto de Requisitos (Cláusula 5.2.6)}
La norma señala en la cláusula 5.2.6 \cite{iso29148} que no basta con tener frases individuales bien redactadas, sino que el grupo completo de requisitos debe funcionar como un todo ordenado y coherente:
\begin{itemize}
    \item \textbf{Completo del conjunto (Set Complete):} Cubre la totalidad del alcance del proyecto sin dejar vacíos funcionales o situaciones sin contemplar.
    \item \textbf{Consistente (Consistent):} Evita contradicciones lógicas entre requerimientos; un requisito no puede pedir lo contrario de otro.
    \item \textbf{Factible del conjunto (Set Feasible):} Todo el grupo de requerimientos puede implementarse al mismo tiempo con los recursos disponibles.
    \item \textbf{Modificable (Modifiable):} El documento está organizado de tal forma que permite hacer cambios futuros en un requisito sin desorganizar todo el informe.
    \item \textbf{Trazable (Traceable):} Asigna un código o identificador único a cada requisito para rastrear su origen y su prueba correspondiente.
\end{itemize}

\subsection{Criterios de Lenguaje y Términos Prohibidos (Cláusula 5.2.7)}
La cláusula 5.2.7 \cite{iso29148} prohíbe taxativamente el empleo de adjetivos, adverbios u oraciones ambiguas que no admitan una métrica objetiva. Términos como \textit{«rápido»}, \textit{«seguro»}, \textit{«fácil de usar»}, \textit{«robusto»}, \textit{«optimizado»}, \textit{«eficiente»}, \textit{«según sea necesario»} o \textit{«en la medida de lo posible»} quedan excluidos por ser indemostrables en un protocolo de pruebas formal. En su lugar, la norma exige la cuantificación mediante indicadores numéricos concretos con métrica, unidad de medida y umbral de aceptación.

\section{Glosario de Términos (Cláusula 3 de ISO/IEC/IEEE 29148)}

A continuación se parafrasean 10 términos clave de la cláusula 3 de la norma \cite{iso29148}, acompañados de un ejemplo contextualizado en la plataforma Mercado Libre:

\begin{enumerate}
    \item \textbf{Requisito (Requirement):} Expresión formal de una capacidad o condición que debe satisfacer un sistema para resolver un problema real o cumplir un contrato. \\
    \textit{Ejemplo en Mercado Libre:} El sistema deberá procesar el pago mediante tarjeta de débito confirmando la transacción con la entidad bancaria.
    \item \textbf{Requisito Funcional (Functional Requirement):} Enunciado que especifica una acción o comportamiento operativo que el software debe ejecutar ante una entrada de datos. \\
    \textit{Ejemplo en Mercado Libre:} El usuario podrá filtrar la lista de resultados seleccionando la categoría «Monitores y Accesorios».
    \item \textbf{Restricción (Constraint):} Exigencia organizativa, legal o técnica que limita el abanico de soluciones de diseño posibles durante el desarrollo. \\
    \textit{Ejemplo en Mercado Libre:} La aplicación deberá cumplir con el estándar PCI-DSS para el manejo seguro de datos bancarios.
    \item \textbf{Requisito de Interfaz Externa (External Interface Requirement):} Especificación de los puntos de conexión y protocolos de intercambio de información con sistemas externos. \\
    \textit{Ejemplo en Mercado Libre:} La app se comunicará con los servicios web del operador logístico para calcular la fecha estimada de entrega.
    \item \textbf{Requisito de Desempeño (Performance Requirement):} Especificación cuantitativa referida a la velocidad, capacidad de procesamiento o rendimiento bajo cargas de trabajo definidas. \\
    \textit{Ejemplo en Mercado Libre:} La búsqueda de productos responderá en un tiempo menor a 1.5 segundos ante solicitudes concurrentes.
    \item \textbf{Atributo (Attribute):} Propiedad o característica cualitativa o cuantitativa asociada a un elemento de requisito (ej. prioridad, estado, autor). \\
    \textit{Ejemplo en Mercado Libre:} El atributo «Prioridad = Alta» asignado al requisito de procesamiento del botón «Pagar y finalizar».
    \item \textbf{Trazabilidad (Traceability):} Capacidad de discernir la relación histórica entre un requisito, su origen de negocio y los componentes de software derivados. \\
    \textit{Ejemplo en Mercado Libre:} Vincular el ID del requisito \texttt{RF-03} directamente con la prueba unitaria del carrito de compras.
    \item \textbf{Verificación (Verification):} Confirmación mediante evidencia objetiva de que el sistema cumple con las especificaciones establecidas en los requisitos. \\
    \textit{Ejemplo en Mercado Libre:} Ejecutar una prueba automatizada para confirmar que el cupón de descuento resta efectivamente US\$ 26.71 al total.
    \item \textbf{Especificación de Requisitos de Software (SRS):} Documento formal estructurado que consolida la totalidad de requerimientos que el software debe satisfacer. \\
    \textit{Ejemplo en Mercado Libre:} El documento técnico que especifica las reglas de negocio para la compra protegida y garantías.
    \item \textbf{Interesado (Stakeholder):} Persona o entidad que afecta o puede verse afectada por las decisiones y resultados del desarrollo del sistema. \\
    \textit{Ejemplo en Mercado Libre:} El comprador final, las tiendas oficiales y el equipo de soporte al cliente.
\end{enumerate}

\section{Evidencias de Campo del Sistema Real}

Para respaldar técnicamente la elicitación y comprobación de los requisitos de Mercado Libre, se incluyen tres capturas de pantalla tomadas directamente de la aplicación móvil (Figura~\ref{fig:capturas}). *(Nota: Asegúrate de guardar tus imágenes con los nombres `cap1.jpeg`, `cap2.jpeg` y `cap3.jpeg` dentro de la carpeta `TA-IND-01/figuras/` de tu repositorio).*

\begin{figure}[h!]
    \centering
    \begin{minipage}[b]{0.31\textwidth}
        \centering
        \includegraphics[width=\textwidth]{captura1.jpeg}
        \caption{Filtros de búsqueda (Monitor).}
        \label{fig:cap1}
    \end{minipage}
    \hfill
    \begin{minipage}[b]{0.31\textwidth}
        \centering
        \includegraphics[width=\textwidth]{captura2.jpeg}
        \caption{Detalle de producto y stock.}
        \label{fig:cap2}
    \end{minipage}
    \hfill
    \begin{minipage}[b]{0.31\textwidth}
        \centering
        \includegraphics[width=\textwidth]{captura3.jpeg}
        \caption{Pantalla de pago y checkout.}
        \label{fig:cap3}
    \end{minipage}
    \caption{Evidencias de pantalla de la app Mercado Libre registradas por el estudiante.}
    \label{fig:capturas}
\end{figure}

\section{Formulación y Clasificación de 15 Requisitos}

La Tabla~\ref{tab:requisitos} resume los 15 requisitos redactados para Mercado Libre, abarcando 5 tipos de la norma ISO/IEC/IEEE 29148 \cite{iso29148} e incluyendo 5 requisitos de calidad cuantificados conforme a la ISO/IEC 25010:2023 \cite{iso25010}.

\small
\begin{longtable}{p{0.08\textwidth} p{0.32\textwidth} p{0.18\textwidth} p{0.18\textwidth} p{0.18\textwidth}}
\caption{Tabla de clasificación y criterios de aceptación de los 15 requisitos.} \label{tab:requisitos} \\
\toprule
\textbf{ID} & \textbf{Enunciado} & \textbf{Tipo (5.2)} & \textbf{ISO 25010 (si aplica)} & \textbf{Criterio de Aceptación Verificable} \\
\midrule
\endhead
\bottomrule
\endfoot
RF01 & El sistema permitirá aplicar filtros múltiples por envío gratis, ofertas y tamaño de pantalla en la búsqueda. & Funcional & N/A & Retornar únicamente productos que coincidan con las casillas activadas. \\
RF02 & El sistema desplegará el distintivo «Más Vendido» y la tienda oficial responsable del producto. & Funcional & N/A & Visualizar el distintivo en la vista detallada si el vendedor supera 1000 ventas. \\
RF03 & El sistema permitirá seleccionar la cantidad deseada mostrando las unidades disponibles (ej. +25 disponibles). & Funcional & N/A & Impedir seleccionar una cantidad mayor al stock registrado en base de datos. \\
RF04 & El sistema calculará el resumen de compra descontando cupones activos (ej. descuento de US\$ 26.71). & Funcional & N/A & Restar el monto del cupón del precio base mostrando el total exacto. \\
RF05 & El sistema registrará la dirección de entrega y actualizará la fecha estimada. & Funcional & N/A & Guardar los datos de envío y recalcular la ventana de llegada de forma automática. \\
RF06 & El sistema permitirá seleccionar «Nueva tarjeta de débito» como método de pago en el checkout. & Funcional & N/A & Presentar el formulario de ingreso de tarjeta y validar los 16 dígitos. \\
RF07 & El sistema incluirá la opción «Compra Protegida» con devolución dentro de los primeros 30 días. & Funcional & N/A & Generar la etiqueta de devolución gratuita si el usuario cancela en el plazo. \\
RC01 & El tiempo de procesamiento al presionar «Ver resultados» en los filtros no superará 1.2 segundos bajo 2000 usuarios concurrentes. & Calidad (Desempeño) & Eficiencia de desempeño - Tiempo & Tiempo de latencia $\le 1.2$ s en servidor bajo prueba de carga. \\
RC02 & Los datos de facturación se cifrarán con algoritmo AES-256 en reposo y TLS 1.3 en tránsito. & Calidad (Seguridad) & Seguridad - Confidencialidad & Cifrado comprobado al 100\% en auditoría de red de paquetes HTTP. \\
RC03 & El flujo completo desde la ficha del producto hasta el checkout se ejecutará en máximo 3 clics en usuarios mayores a 60 años. & Calidad (Usabilidad) & Usabilidad - Operabilidad & Prueba de usabilidad con 90\% de éxito en $\le 3$ interacciones. \\
RC04 & La aplicación mantendrá una disponibilidad operativa de 99.9\% medido en ciclo mensual 24/7. & Calidad (Fiabilidad) & Fiabilidad - Disponibilidad & Registro de caídas $\le 43.8$ minutos de interrupción no planificada al mes. \\
RC05 & La integración con la API del proveedor de facturación devolverá respuesta exitosa en $\ge 99.8\%$ de intentos. & Calidad (Compatibilidad) & Compatibilidad - Interoperabilidad & Tasa de error en llamadas a la API externa $\le 0.2\%$. \\
RES1 & La aplicación móvil deberá ejecutarse en dispositivos Android 10 o versiones posteriores. & Restricción & N/A & Instalación y ejecución correcta en compilación con API Level 29+. \\
INT1 & El módulo de ubicación consultará el código postal del usuario a través del servicio GPS del teléfono. & Interfaz externa & N/A & Recepción de coordenadas geográficas en formato decimal con precisión de 10m. \\
RES2 & La interfaz responderá al estándar de colores azul corporativo y amarillo sin alterar botones clave. & Restricción & N/A & Inspección de hoja de estilos confirmando códigos hexadecimales oficiales. \\
\end{longtable}
\normalsize

\section{Autoevaluación y Corrección de Requisitos}

En apego a la norma ISO/IEC/IEEE 29148 (cláusulas 5.2.5 y 5.2.7) \cite{iso29148}, se evaluaron las 15 formulaciones iniciales identificando fallos de ambigüedad, falta de singularidad o falta de umbrales cuantitativos. La Tabla~\ref{tab:autoevaluacion} ilustra las revisiones aplicadas.

\begin{table}[h!]
\small
\centering
\begin{tabularx}{\textwidth}{l X X X}
\toprule
\textbf{ID} & \textbf{Versión Anterior (Con deficiencias)} & \textbf{Fallo Detectado (5.2.5 / 5.2.7)} & \textbf{Versión Corregida (Normativa)} \\
\midrule
RC01 & La búsqueda con filtros debe ser rápida cuando el usuario navegue. & Contiene el término prohibido «rápida» (5.2.7) y no es verificable. & El tiempo de respuesta de la búsqueda con filtros no superará los 1.2 segundos bajo 2000 usuarios concurrentes. \\
RC02 & Los datos de la cédula del usuario deben estar bien protegidos en el sistema. & Ambiguo y subjetivo («bien protegidos»); viola 5.2.5 (Verificable). & Los datos de facturación se cifrarán obligatoriamente mediante AES-256 en reposo y TLS 1.3 en tránsito. \\
RC03 & El sistema debe ser fácil de usar y comprar rápido en pocos pasos. & Múltiples adjetivos prohibidos («fácil», «rápido») y falta singularidad. & El flujo de compra directa se ejecutará en un máximo de 3 clics desde la ficha del producto. \\
RF04 & El sistema aplicará descuentos y calculará los envíos y mostrará el total. & Viola la propiedad de Singularidad (5.2.5) al agrupar tres acciones distintas. & El sistema calculará el total descontando cupones activos (ej. US\$ 26.71) y sumando el costo de envío. \\
RC04 & La app no debe caerse casi nunca durante el año de funcionamiento. & Término ambiguo («casi nunca»); viola la característica de Inambiguo. & La aplicación mantendrá una disponibilidad operativa del 99.9\% medida en ciclos mensuales. \\
\bottomrule
\end{tabularx}
\caption{Tabla de autoevaluación: versión defectuosa vs. versión corregida.}
\label{tab:autoevaluacion}
\end{table}

\section{Reflexión Final}

Durante la formulación y cuantificación de requisitos para la plataforma Mercado Libre, el tipo de requisito que representó la mayor complejidad conceptual y técnica fue el de \textbf{Calidad en Uso y Usabilidad cuantificada (RC03)}, bajo el modelo ISO/IEC 25010:2023 \cite{iso25010}. 

La dificultad principal radica en la tendencia habitual a redactar la usabilidad mediante adjetivos cualitativos tradicionales («intuitivo», «fácil de usar», «amigable»). Adaptar estas necesidades a enunciados rigurosos según la cláusula 5.2.7 de la ISO/IEC/IEEE 29148 exige definir métricas cuantitativas observables ---tales como la cantidad máxima de clics, la tasa de éxito de la tarea o la tolerancia a errores del usuario--- sin invadir ni condicionar de forma prematura las decisiones de diseño de la interfaz gráfica. Asimismo, el equilibrio entre fijar un umbral de rendimiento realista (como latencias de 1.2s en la búsqueda) y no asumir arquitecturas técnicas inexistentes requiere una lectura minuciosa del contexto operativo real del software.

\section{Declaración de Uso de Inteligencia Artificial}
\begin{itemize}
    \item \textbf{Herramienta y Versión:} Gemini.
    \item \textbf{Propósito:} Corrección ortográfica del texto y apoyo en la estructuración de sintaxis de tablas de alta densidad en \LaTeX.
    \item \textbf{Verificación del estudiante:} Se revisaron manualmente los enunciados, asegurando que la clasificación corresponda a las cláusulas reales de la norma ISO/IEC/IEEE 29148:2018 \cite{iso29148} y que los datos informativos pertenezcan a las capturas reales obtenidas del sistema Mercado Libre.
\end{itemize}

\begin{thebibliography}{99}
\bibitem{iso29148} ISO/IEC/IEEE, \emph{Systems and software engineering -- Life cycle processes -- Requirements engineering}, ISO/IEC/IEEE Standard 29148:2018, 2018.
\bibitem{iso25010} ISO/IEC, \emph{Systems and software engineering -- Systems and software quality requirements and evaluation (SQuaRE) -- Product quality model}, ISO/IEC Standard 25010:2023, 2023.
\bibitem{pohl2025} K. Pohl, \emph{Requirements Engineering: Fundamentals, Principles, and Techniques}, 2nd ed. Berlin, Heidelberg, Germany: Springer, 2025.
\bibitem{swebok2024} H. Washizaki, Ed., \emph{Guide to the Software Engineering Body of Knowledge (SWEBOK Guide)}, Version 4.0. Los Alamitos, CA, USA: IEEE Computer Society, 2024.
\bibitem{lopdp2021} Asamblea Nacional del Ecuador, \emph{Ley Orgánica de Protección de Datos Personales}, Registro Oficial, Quinto Suplemento No. 459, 26 de mayo de 2021.
\bibitem{acm2018} Association for Computing Machinery, \emph{ACM Code of Ethics and Professional Conduct}, https://www.acm.org/code-of-ethics, 2018.
\end{thebibliography}

\end{document}
